
\chapter*{Περίληψη}

Τα μικρά γλωσσικά μοντέλα μπορούν να εκτελούνται σε κινητά τηλέφωνα, αλλά η διατήρηση
μακρών αλληλεπιδράσεων επιβαρύνει το περιορισμένο παράθυρο συμφραζομένων, την μπαταρία, τη
μνήμη και τα θερμικά τους όρια. Οι μελέτες σε φορητές συσκευές που εξετάσαμε δεν συγκρίνουν
αναπαραστάσεις μνήμης για ένα ιστορικό συνομιλίας ως προς την ποιότητα των απαντήσεων και την
ενέργεια στην ίδια συσκευή. Η παρούσα διπλωματική εργασία εξετάζει πώς οι στρατηγικές
μνήμης επηρεάζουν την ποιότητα των απαντήσεων, ποιους συμβιβασμούς πόρων δημιουργούν σε ένα
κινητό τηλέφωνο και πού η μνήμη με μορφή γράφου γνώσης επιτυγχάνει ή αποτυγχάνει.

Μια κοινή υλοποίηση συγκρίνει το πλήρες ιστορικό με περικοπή, τη διανυσματική ανάκτηση και
γράφους που εξάγονται από μικρά γλωσσικά μοντέλα. Η κύρια μελέτη εκτελείται σε ένα Samsung
Galaxy S23 Ultra και μετρά την ακρίβεια, την ενέργεια, τον χρόνο απόκρισης, τη χρήση μνήμης
και τη θερμοκρασία, αναφέροντας την ακρίβεια και την ενέργεια χωριστά και ως σωστές
απαντήσεις ανά κιλοτζάουλ (kJ). Ένας σταθμός εργασίας εκτελεί τις ίδιες μονάδες με
περισσότερες ερωτήσεις και μοντέλα Gemma και Qwen στα LoCoMo και LongMemEval, και εξετάζει το βάθος ανάκτησης, τις διορθώσεις των
γράφων και την επαναφορά του αρχικού κειμένου.

Στο τηλέφωνο, η διανυσματική ανάκτηση απάντησε σωστά σε περισσότερες ερωτήσεις από το πλήρες
ιστορικό με περικοπή και κατανάλωσε λιγότερη ενέργεια ανά σωστή απάντηση. Στον σταθμό
εργασίας, το πλήρες ιστορικό με μεγάλο παράθυρο συμφραζομένων είχε υψηλότερες σημειακές
εκτιμήσεις ακρίβειας από τη διανυσματική ανάκτηση στο LoCoMo και χαμηλότερες στο LongMemEval-S,
με πολλαπλάσιο μήκος εισόδου. Σε μία πιλοτική
κατασκευή στο τηλέφωνο, η δημιουργία γράφου κόστισε σημαντικά περισσότερη ενέργεια από τη
δημιουργία του διανυσματικού ευρετηρίου. Η ακρίβεια της μνήμης γράφου μετρήθηκε στον
σταθμό εργασίας. Εκεί, η μνήμη γράφου περιορίστηκε από αποτυχίες εξαγωγής και από τον τρόπο με
τον οποίο τα στοιχεία έφταναν στο μοντέλο, ενώ η επαναφορά του αρχικού κειμένου πίσω από τα
ανακτημένα γεγονότα βελτίωσε τις απαντήσεις των γράφων για τα περισσότερα μοντέλα, χωρίς να
φτάσει τη διανυσματική ανάκτηση. Η μέθοδος Επιλεκτικής Επέκτασης Πηγών (Selective
Provenance Expansion, SPE), που αναπτύχθηκε για να επιλέγει ανάμεσα σε συμπαγή γεγονότα και
αρχικό κείμενο υπό όριο λεκτικών μονάδων (tokens), δεν ανέκτησε με συνέπεια αυτό το όφελος.

Τα ευρήματα συνηγορούν στην αξιολόγηση της μνήμης τόσο ως προς τα στοιχεία που παρέχει όσο
και ως προς τους πόρους που απαιτεί η κατασκευή και η χρήση της. Περιορίζονται από τη μία
συσκευή, το μικρό δείγμα στο τηλέφωνο και την αυτόματη βαθμολόγηση, και αφορούν
ερωταποκρίσεις και όχι ολοκληρωμένες ροές εργασίας κινητών πρακτόρων.

\medskip
\noindent\textbf{Λέξεις-κλειδιά:} μικρά γλωσσικά μοντέλα, μακροπρόθεσμη μνήμη,
παραγωγή επαυξημένη με ανάκτηση, γράφοι γνώσης, τοπική εκτέλεση μοντέλων,
ενεργειακή αποδοτικότητα.
